
\documentclass{article}
\usepackage{colortbl}
\usepackage{makecell}
\usepackage{multirow}
\usepackage{supertabular}

\begin{document}

\newcounter{utterance}

\centering \large Interaction Transcript for game `cladder', experiment `full\_v1.5\_default', episode 2677 with refineonly{-}Qwen3.5{-}9B.
\vspace{24pt}

{ \footnotesize  \setcounter{utterance}{1}
\setlength{\tabcolsep}{0pt}
\begin{supertabular}{c@{$\;$}|p{.15\linewidth}@{}p{.15\linewidth}p{.15\linewidth}p{.15\linewidth}p{.15\linewidth}p{.15\linewidth}}
    \# & \multicolumn{2}{c}{Player} && \multicolumn{2}{c}{Game Master} \\
    \hline

    \theutterance \stepcounter{utterance}  
    & & & \multicolumn{4}{p{0.6\linewidth}}{
        \cellcolor[rgb]{0.9,0.9,0.9}{
            \makecell[{{p{\linewidth}}}]{
                \texttt{\tiny{[P1$\langle$GM]}}
                \texttt{You are an expert in causal inference. The following question is not a typical commonsense query, but rather a meticulously designed question created by a professor specializing in causal inference, intended to assess the students' mastery of the course content.} \\
\texttt{Imagine a self{-}contained, hypothetical world with only the following conditions, and without any unmentioned factors or causal relationships: Gene has a direct effect on smoking and lung cancer. Pollution has a direct effect on smoking and lung cancer. Smoking has a direct effect on lung cancer. Pollution is unobserved. Method 1: We look at how gene correlates with lung cancer case by case according to smoking. Method 2: We look directly at how gene correlates with lung cancer in general. To understand how gene affects lung cancer, is it more correct to use the Method 1 than Method 2?} \\
\texttt{Start your answer with 'Yes' or 'No', followed by additional reasoning or evidence to support your explanation.} \\
            }
        }
    }
     \\ \\

    \theutterance \stepcounter{utterance}  
    & & \multicolumn{4}{p{0.6\linewidth}}{
        \cellcolor[rgb]{0.95,0.95,0.95}{
            \makecell[{{p{\linewidth}}}]{
                \texttt{\tiny{[P1$\leftrightarrow$P]}}
                \texttt{{[}DRAFT{]}} \\
\texttt{No} \\
\\ 
\texttt{**Reasoning:**} \\
\\ 
\texttt{In this hypothetical world, the causal structure can be represented as follows:} \\
\texttt{1.  **Gene** \$\textbackslash{}rightarrow\$ **Smoking**} \\
\texttt{2.  **Gene** \$\textbackslash{}rightarrow\$ **Lung Cancer**} \\
\texttt{3.  **Pollution** \$\textbackslash{}rightarrow\$ **Smoking**} \\
\texttt{4.  **Pollution** \$\textbackslash{}rightarrow\$ **Lung Cancer**} \\
\texttt{5.  **Smoking** \$\textbackslash{}rightarrow\$ **Lung Cancer**} \\
\\ 
\texttt{Here, **Smoking** acts as a **mediator** through which Gene influences Lung Cancer, but it also acts as a **collider** (or part of a path) involving Pollution. More critically, because **Pollution** is an unobserved confounder that affects both **Smoking** and **Lung Cancer**, and **Gene** affects **Smoking**, conditioning on **Smoking** (as done in Method 1) opens a non{-}causal backdoor path between **Gene** and **Lung Cancer**.} \\
\\ 
\texttt{Let's trace the paths:} \\
\texttt{*   **Direct Path:** Gene \$\textbackslash{}rightarrow\$ Lung Cancer (This is the causal effect we want to estimate).} \\
\texttt{*   **Induced Path via Conditioning on Smoking:**} \\
\texttt{    *   Gene \$\textbackslash{}rightarrow\$ Smoking} \\
\texttt{    *   Smoking \$\textbackslash{}leftarrow\$ Pollution \$\textbackslash{}rightarrow\$ Lung Cancer} \\
\texttt{    *   By conditioning on **Smoking**, we induce a correlation between **Gene** and **Pollution** (collider bias/fork opening). Since **Pollution** causes **Lung Cancer**, this creates a spurious association between **Gene** and **Lung Cancer** that is not causal.} \\
\\ 
\texttt{**Analysis of the Methods:**} \\
\texttt{*   **Method 2 (Unconditional correlation):** This estimates the total effect of Gene on Lung Cancer, which includes both the direct effect and the indirect effect mediated through Smoking. While this does not isolate the *direct* effect, it does not introduce the specific bias caused by conditioning on the mediator in the presence of an unobserved confounder affecting the mediator. However, if the goal is strictly the *direct* effect, Method 2 is biased by the indirect path.} \\
\texttt{*   **Method 1 (Conditional correlation):** This attempts to isolate the direct effect by holding Smoking constant. However, because **Pollution** is unobserved and affects both Smoking and Lung Cancer, conditioning on Smoking creates a spurious association between Gene and Pollution. Consequently, the correlation observed in Method 1 will be a mixture of the true direct effect and a large spurious effect induced by the unobserved confounder Pollution. This makes the estimate in Method 1 highly biased and unreliable for inferring the causal mechanism.} \\
\\ 
\texttt{Therefore, in the presence of an unobserved confounder (Pollution) that affects the mediator (Smoking), conditioning on the mediator (Method 1) introduces severe bias (M{-}bias or similar collider bias structures depending on the} \\
            }
        }
    }
    & \\ \\

    \theutterance \stepcounter{utterance}  
    & \multicolumn{4}{p{0.6\linewidth}}{
        \cellcolor[rgb]{0.561,0.737,0.561}{
            \makecell[{{p{\linewidth}}}]{
                \texttt{\tiny{[P1$\rangle$GM]}}
                \texttt{No} \\
\\ 
\texttt{**Reasoning:**} \\
\\ 
\texttt{In this hypothetical world, the causal structure is defined by the following directed edges:} \\
\texttt{1.  **Gene** \$\textbackslash{}rightarrow\$ **Smoking**} \\
\texttt{2.  **Gene** \$\textbackslash{}rightarrow\$ **Lung Cancer**} \\
\texttt{3.  **Pollution** \$\textbackslash{}rightarrow\$ **Smoking**} \\
\texttt{4.  **Pollution** \$\textbackslash{}rightarrow\$ **Lung Cancer**} \\
\texttt{5.  **Smoking** \$\textbackslash{}rightarrow\$ **Lung Cancer**} \\
\\ 
\texttt{Here, **Smoking** acts as a **mediator** for the effect of Gene on Lung Cancer. However, because **Pollution** is an unobserved variable that affects both **Smoking** and **Lung Cancer**, it creates a confounding relationship between the mediator and the outcome.} \\
\\ 
\texttt{Let's analyze the paths when conditioning on **Smoking** (as done in Method 1):} \\
\texttt{*   Conditioning on **Smoking** (a node with two incoming arrows: from Gene and from Pollution) turns it into a **collider**.} \\
\texttt{*   This conditioning opens the non{-}causal backdoor path: **Gene** \$\textbackslash{}rightarrow\$ **Smoking** \$\textbackslash{}leftarrow\$ **Pollution** \$\textbackslash{}rightarrow\$ **Lung Cancer**.} \\
\texttt{*   By opening this path, we induce a spurious correlation between **Gene** and **Pollution**. Since **Pollution** is a cause of **Lung Cancer**, this creates a biased estimate of the relationship between **Gene** and **Lung Cancer** that conflates the true direct effect with the confounding effect of Pollution. This is a classic case of **M{-}bias** (or collider bias).} \\
\\ 
\texttt{**Analysis of the Methods:**} \\
\texttt{*   **Method 1 (Conditional correlation):** This method conditions on the mediator (**Smoking**) to attempt to isolate the direct effect. However, due to the unobserved confounder (**Pollution**) affecting the mediator, this conditioning opens a backdoor path, rendering the estimate highly biased and invalid for causal inference.} \\
\texttt{*   **Method 2 (Unconditional correlation):** This method estimates the total effect of **Gene** on **Lung Cancer** (direct + indirect via Smoking). While it does not isolate the *direct* effect specifically, it does not suffer from the collider bias introduced by conditioning on the mediator in the presence of an unobserved confounder. In causal inference, an unbiased estimate of the total effect is generally preferred over a severely biased estimate of the direct effect when the latter cannot be identified without further assumptions or data.} \\
\\ 
\texttt{Therefore, because Method 1 introduces severe bias due to conditioning on a collider in the presence of an unobserved confounder, it is **not** more correct than Method 2 for understanding the causal influence of Gene on Lung Cancer in this specific setup. Method 2 provides a valid (though total) effect estimate, whereas Method 1 provides a fundamentally flawed one.} \\
            }
        }
    }
    & & \\ \\

    \theutterance \stepcounter{utterance}  
    & & & \multicolumn{2}{p{0.3\linewidth}}{
        \cellcolor[rgb]{0.9,0.9,0.9}{
            \makecell[{{p{\linewidth}}}]{
                \texttt{\tiny{[GM$|$GM]}}
                \texttt{No} \\
            }
        }
    }
    & & \\ \\

    \theutterance \stepcounter{utterance}  
    & & & \multicolumn{2}{p{0.3\linewidth}}{
        \cellcolor[rgb]{0.9,0.9,0.9}{
            \makecell[{{p{\linewidth}}}]{
                \texttt{\tiny{[GM$|$GM]}}
                \texttt{no} \\
            }
        }
    }
    & & \\ \\

    \theutterance \stepcounter{utterance}  
    & & & \multicolumn{2}{p{0.3\linewidth}}{
        \cellcolor[rgb]{0.9,0.9,0.9}{
            \makecell[{{p{\linewidth}}}]{
                \texttt{\tiny{[GM$|$GM]}}
                \texttt{game\_result = WIN} \\
            }
        }
    }
    & & \\ \\

\end{supertabular}
}

\end{document}
